%==============================================================================
% Flow-Matching Augmentation for Rare Attack Detection in NIDS
% Introduction section draft.
%
% Compile:  pdflatex main && bibtex main && pdflatex main && pdflatex main
%
% Numbers that depend on model runs are defined as macros in the block below.
% They will change when experiments are re-run; edit them in one place rather
% than hunting through the prose. Values marked FIXED are properties of the
% datasets and will not change.
%==============================================================================

\documentclass[journal]{IEEEtran}

\usepackage{amsmath}
\usepackage{booktabs}
\usepackage{graphicx}
\usepackage{url}
\usepackage[hidelinks]{hyperref}

%------------------------------------------------------------------------------
% Volatile numbers -- update after each experimental re-run.
%------------------------------------------------------------------------------
% Values below are means over 5 seeds on the official KDDTrain+/KDDTest+ split,
% from results/nsl_kdd_resampling_summary.csv. XGBoost uses subsample=0.8 and
% colsample_bytree=0.8 so the seed perturbs the model as well as the sampler;
% with both at 1.0 the classifier is deterministic and every seed returns an
% identical result, reporting a standard deviation of zero that measures nothing.
\newcommand{\baseAcc}{0.781}          % overall accuracy, no augmentation
\newcommand{\baseMacroF}{0.569}       % macro-F1, no augmentation
\newcommand{\baseRtwolRec}{0.107}     % R2L recall, no augmentation
\newcommand{\baseRtwolPrec}{0.974}    % R2L precision, no augmentation
\newcommand{\baseUtwrRec}{0.131}      % U2R recall, no augmentation
\newcommand{\baseRtwolFone}{0.192}    % R2L F1, no augmentation (sd 0.025)
\newcommand{\baseUtwrFone}{0.219}     % U2R F1, no augmentation (sd 0.018)
\newcommand{\bestSmoteUtwr}{0.428}    % best U2R F1 among classical arms (SMOTE, sd 0.027)
\newcommand{\bestAdasynRtwol}{0.251}  % best R2L F1 among classical arms (ADASYN, sd 0.015)
\newcommand{\bestMacroF}{0.627}       % best macro-F1 among classical arms (SMOTE)

% Fidelity-versus-utility, measured directly (see Discussion). ADASYN's synthetic
% rows are perfectly separable from real ones yet yield the best R2L F1 of any arm.
\newcommand{\smoteDetectAUC}{0.535}   % SMOTE, R2L: real-vs-synthetic detection AUC
\newcommand{\adasynDetectAUC}{1.000}  % ADASYN, U2R: detection AUC

%------------------------------------------------------------------------------
% Dataset facts (FIXED -- computed from the public files, will not change)
%------------------------------------------------------------------------------
\newcommand{\trainRows}{125{,}973}
\newcommand{\numnormal}{67{,}343}     % NSL-KDD majority (normal) class, train split
\newcommand{\utwrTrain}{52}
\newcommand{\utwrRatio}{1{:}1{,}295}
\newcommand{\rtwolTrain}{995}
\newcommand{\warezTrain}{890}
\newcommand{\rtwolEffective}{105}
\newcommand{\rtwolEffectiveRatio}{1{:}641}

\begin{document}

\title{When Augmentation Hurts: A Per-Class Evaluation of\\
Generative and Classical Resampling for Rare-Attack Detection}

\author{Prayag Sharma% <-- confirm preferred form of name and affiliation
\thanks{Manuscript draft.}}

\maketitle

\begin{abstract}
Class imbalance is the standard explanation for why intrusion detection systems miss
rare attacks, and synthetic augmentation is the standard remedy. Existing evaluations
of that remedy, however, use binary labels, aggregate metrics and a single downstream
classifier --- a setting in which the rare categories of interest are not separately
observable. We report a per-class evaluation of eight augmentation methods spanning
classical resampling, generative adversarial, diffusion and flow-matching families,
across three benchmark datasets, nine rare attack classes, five seeds and two
classifier families, with Holm--Bonferroni-corrected paired tests and bootstrap
confidence intervals throughout. Augmentation improves rare-class F1 in eleven cases,
by up to $85.6\%$ relative ($p=0.002$, $d=4.68$). It also \emph{degrades} detection
significantly in fourteen; the largest single effect we measure is a harm, SMOTE
reducing UNSW-NB15 Shellcode F1 by $0.139$ at $d=-9.59$. Five different methods are
optimal across the nine classes, and changing only the downstream classifier changes
the best-performing method for all nine. Four widely used sample-quality measures ---
discriminator detectability, nearest-neighbour distance, marginal divergence and
attribution similarity --- none predict downstream utility, and learned generators
fail silently on approximately one run in five where classical resamplers never do.
We conclude that augmentation-method selection for rare-attack detection is a
per-class, per-classifier empirical question, and that applying a default resampler
is as likely to harm rare-attack detection as to help it. Code and all result files
are publicly available.
\end{abstract}

\begin{IEEEkeywords}
Intrusion detection, class imbalance, generative models, flow matching,
synthetic data, rare attack detection.
\end{IEEEkeywords}

%==============================================================================
\section{Introduction}
%==============================================================================

\IEEEPARstart{N}{etwork} intrusion detection systems (NIDS) built on supervised
machine learning are now standard components of enterprise security monitoring,
and the published accuracy figures attached to them are remarkable. Models
evaluated on the common benchmarks routinely report overall accuracy above
99\,\% \cite{alsubaei2025smart, yang2025diffids}. Read alongside the operational
reality that intrusions continue to succeed, those figures invite a question
that the aggregate metric cannot answer: accuracy on \emph{what}?

The answer is uncomfortable. Network traffic is dominated by benign flows, and
attack traffic is dominated by a small number of high-volume attack families.
In the NSL-KDD benchmark \cite{tavallaee2009detailed}, the user-to-root (U2R)
category is represented by \utwrTrain{} of \trainRows{} training records, a ratio
of roughly \utwrRatio{}. A classifier that never predicts U2R at all still scores
above 99.9\,\% on that class by the accuracy metric. The incentive structure of
the loss function makes this outcome not merely possible but expected, and the
headline number conceals it completely.

This matters because rarity and severity are correlated in the wrong direction.
U2R attacks describe privilege escalation, in which an adversary already holding
a normal user account acquires administrative control. Remote-to-local (R2L)
attacks describe unauthorised access obtained from outside the network. Both
categories sit at the point where an intrusion becomes materially damaging, and
both are precisely the categories that vanish into the rounding error of an
aggregate score. Denial-of-service traffic, by contrast, is abundant, easily
learned, and comparatively survivable.

The problem is documented but under-diagnosed. Alsubaei \cite{alsubaei2025smart}
reports above 99\,\% overall accuracy on NSL-KDD, UNSW-NB15
\cite{moustafa2015unsw}, and CICIDS2017 \cite{sharafaldin2018toward} after
systematic hyperparameter optimisation, while a rare class in the same
evaluation remains at an F1-score near 0.15. Careful tuning of a strong model,
in other words, does not touch the failure. Our own unaugmented baseline
reproduces the pattern: on the official NSL-KDD split we obtain overall accuracy
of \baseAcc{} against a macro-averaged F1 of \baseMacroF{}, with R2L recall of
\baseRtwolRec{} and U2R recall of \baseUtwrRec{}. The R2L precision of
\baseRtwolPrec{} is instructive. The classifier is almost never wrong when it
does flag an R2L attack; it simply declines to flag them.

\subsection{Existing responses and their limits}

Two families of remedy dominate the literature. The first resamples the training
distribution. SMOTE \cite{chawla2002smote} synthesises minority examples by
interpolating between a sample and its nearest same-class neighbours, and ADASYN
\cite{he2008adasyn} biases that interpolation toward regions where minority
samples are locally outnumbered. Both are inexpensive, both are two decades old,
and both remain the practical default. Their constraint is structural: an
interpolation of existing points cannot introduce information absent from those
points, and on data whose class boundaries are non-convex the midpoint of two
minority samples may fall inside a majority region, injecting mislabelled
training data.

The second family learns a generative model of the minority distribution and
samples from it. Conditional tabular GANs \cite{xu2019modeling} were an early
instantiation; denoising diffusion probabilistic models \cite{ho2020denoising},
adapted to tabular data by TabDDPM \cite{kotelnikov2023tabddpm}, are the current
preference. Applications to intrusion detection have followed quickly, with
diffusion-based augmentation reported across the standard benchmarks
\cite{yang2025diffids}. The reported gains are large.

They are also difficult to interpret, for three reasons.

First, the comparisons are incomplete. Diff-IDS \cite{yang2025diffids} evaluates
against deep architectures --- DNN, CNN, RNN, LSTM --- and reports 99.93\,\%
accuracy, but does not compare against SMOTE. A generative model that
outperforms an unaugmented classifier has demonstrated that augmentation helps;
it has not demonstrated that \emph{generative} augmentation helps, which is the
claim being made. Establishing the latter requires the classical resampler as a
control, and that control is frequently absent.

Second, the evaluation frequently omits the information needed to assess a
rare-class claim. Per-class support is often unreported, so a reader cannot
distinguish a stable estimate from one computed over a handful of test samples.
Single-run results are common, which conflates method effect with seed variance.
Both omissions are consequential precisely in the regime the papers claim to
address.

Third, and most fundamentally, the field has not established \emph{why}
augmentation succeeds when it succeeds. The implicit argument for preferring a
generative model over interpolation is that its samples are more realistic, and
that realism should translate into better detection. We find that it does not.
Training a discriminator to separate synthetic from real minority samples, SMOTE
is near-indistinguishable from real data on NSL-KDD R2L (AUC
\smoteDetectAUC{}), while ADASYN is perfectly separable on U2R (AUC
\adasynDetectAUC{}). Yet ADASYN yields the highest R2L F1 of any method we test.
Sample fidelity and downstream utility come apart, and a synthetic-data
evaluation that reports only the former can rank methods in the wrong order.

\subsection{This work}

We address the rare-attack problem with flow matching
\cite{lipman2023flow}, a continuous-time generative framework that learns a
velocity field transporting a simple prior to the data distribution. Flow
matching belongs to the same family as diffusion but differs in mechanism: where
diffusion learns to invert a stochastic corruption process over many discrete
steps, flow matching regresses a deterministic velocity along a prescribed
probability path. The practical consequences are fewer integration steps at
sampling time, a simpler training objective, and no noise schedule to tune. To
our knowledge, flow matching has not previously been applied to class imbalance
in network intrusion detection.

We pair the method with an evaluation protocol designed to answer the questions
the existing comparisons leave open. Every augmentation arm --- no augmentation,
random oversampling, SMOTE, ADASYN, CTGAN, TabDDPM, and flow matching --- passes
through an identical pipeline, over multiple seeds, with per-class precision,
recall, F1, and test support reported throughout. Synthetic data is generated
from the training split alone and added to the training split alone; the test
split is transformed using statistics fitted on training data and is evaluated
once.

We also report a structural property of NSL-KDD that constrains what
augmentation can achieve on that benchmark and, we argue, has been widely
misread. Of the \rtwolTrain{} R2L records in the training split, \warezTrain{}
belong to a single attack type, \texttt{warezclient}, which does not appear in
the test split at all. Conversely, the largest R2L attack types in the test
split include several with no training representation whatsoever. The usable R2L
training signal is therefore closer to \rtwolEffective{} records than to
\rtwolTrain{}, an effective imbalance of approximately
\rtwolEffectiveRatio{} rather than the \(1{:}68\) implied by the class counts. A
generator fitted to the R2L class as a whole learns predominantly to reproduce
\texttt{warezclient}. This reframes the widely cited R2L failure on this
benchmark as substantially a problem of distribution shift rather than of class
balance, and it motivates generating at the level of individual attack types
rather than coarse categories.

\subsection{Contributions}

\begin{enumerate}
\item We introduce flow matching as a minority-class augmenter for network
intrusion detection, and describe a tabular formulation that handles mixed
categorical and continuous features.

\item We provide the controlled comparison the recent generative-augmentation
literature omits, evaluating flow matching against classical resamplers and
against GAN- and diffusion-based generators under an identical pipeline, with
per-class metrics, reported support, multiple seeds, and significance testing.

\item We assess synthetic-sample quality directly rather than inferring it from
downstream accuracy, using train-on-synthetic/test-on-real transfer, a
real-versus-synthetic discriminability test, and nearest-neighbour distance to
detect memorisation.

\item We show that the rare-class failure on NSL-KDD is driven substantially by
train/test attack-type disjointness rather than by class imbalance, quantify the
resulting gap between nominal and effective minority sample counts, and identify
the conditions under which augmentation can and cannot help.

\item We release the full pipeline, including multi-class preprocessing and
label mappings, with instructions sufficient to reproduce every reported table.
\end{enumerate}

\subsection{Organisation}

Section~\ref{sec:related} reviews imbalance handling in intrusion detection and
generative modelling for tabular data. Section~\ref{sec:method} formalises the
flow-matching augmenter. Section~\ref{sec:setup} details datasets, baselines,
and the evaluation protocol. Section~\ref{sec:results} reports results,
ablations, and synthetic-quality diagnostics. Section~\ref{sec:discussion}
examines where augmentation helps and where it cannot, and
Section~\ref{sec:conclusion} concludes.

%==============================================================================
\section{Related Work}\label{sec:related}
%==============================================================================
\subsection{Resampling for imbalanced intrusion detection}

The dominant response to rare-attack detection is to rebalance the training
distribution. SMOTE~\cite{chawla2002smote} synthesises minority records by
interpolating between a record and its nearest same-class neighbours;
ADASYN~\cite{he2008adasyn} biases that sampling toward records whose neighbourhoods
are dominated by other classes. Both are two decades old, both are available in one
line from standard libraries, and both appear in a large fraction of intrusion
detection papers reporting improved rare-class performance.

Their limitations are known in the general machine learning literature. Interpolation
between neighbours cannot leave the convex hull of the minority records, so it adds
density rather than information, and Elor and Averbuch-Elor~\cite{elor2022smote}
report that balancing offers little benefit once a sufficiently strong classifier is
used. What has not been reported, and what we quantify here, is that these methods
can reduce rare-class F1 substantially \emph{below} the unaugmented baseline, with
corrected significance and large effect sizes, on the very classes they are applied
to protect.

\subsection{Generative models for tabular data}

Interpolation has increasingly been replaced by learned generators.
CTGAN~\cite{xu2019modeling} adapts adversarial training to tabular data through
mode-specific normalisation of continuous columns and a conditional sampler that
addresses category imbalance. Denoising diffusion~\cite{ho2020denoising} was adapted
to the tabular setting by TabDDPM~\cite{kotelnikov2023tabddpm}, which applies
Gaussian diffusion to numerical fields and multinomial diffusion to categorical and
binary ones --- an explicit recognition that mixed-type records cannot be treated as
points in a Euclidean space.

Flow matching~\cite{lipman2023flow} trains a velocity field by regression against a
prescribed probability path, avoiding both adversarial dynamics and long ancestral
sampling chains. It is widely used in continuous domains but, to our knowledge, has
not previously been applied to intrusion detection augmentation. Its practical claim
over diffusion --- comparable sample quality at far fewer network evaluations --- is
one we test directly.

\subsection{Generative augmentation for intrusion detection}

Applications of generative models to intrusion datasets are numerous, and several
recent works train a separate generator per attack category rather than a single
conditional model. Diffusion has been applied to this problem specifically, including
work using the benchmarks studied here~\cite{yang2025diffids}. A separate line of
work examines how far intrusion detection models transfer between datasets at
all~\cite{cantone2024machine}, finding substantial degradation; our observation that
the three benchmarks disagree on the direction of the SMOTE effect is a
complementary result at the level of the augmentation method rather than the
detector.

Three features recur across this literature and together define the gap this paper
addresses. Evaluation is typically performed with a single downstream classifier.
Results are typically reported as aggregate accuracy or macro-F1 rather than per
class. And configurations are typically run once, without repetition or significance
testing, so that differences of a few hundredths of an F1 point are presented as
improvements. Our measurements indicate that all three practices are consequential: a
second classifier gives a different answer on every rare class we examine, aggregate
metrics conceal effects of opposite sign occurring simultaneously within one dataset,
and roughly one learned-generator run in five fails silently in a way that a single
run would either miss entirely or report as a result.

\subsection{Evaluating synthetic data}

The closest recent work to ours is the architectural selection framework
of~\cite{architectural2026selection}, which benchmarks twelve generative
architectures on NSL-KDD and CIC-IDS2017, measuring structural fidelity, correlation
preservation and distributional agreement before assessing utility, over twenty
repetitions with paired significance tests. It concludes that GAN-based architectures
achieve the most stable fidelity--utility balance and that diffusion models incur
prohibitive cost at scale.

That study and this one ask different questions. It evaluates \emph{replacement} ---
training on synthetic data and testing on real --- under binary attack/benign labels,
and reports aggregate F1 and accuracy. We evaluate \emph{augmentation}, where
synthetic records supplement rather than replace real ones, under multi-class labels,
and report per-class results for nine individually named rare attack categories. The
distinction matters for two reasons. Under a binary label the rare categories of
interest are absorbed into a single attack class, so effects specific to them are not
separately observable. And a method that produces a usable stand-alone training set
need not be the method that best supplements an existing one. We take that work as
the reference point for fidelity benchmarking and address the per-class utility
question it leaves open.

Beyond intrusion detection, a broader literature evaluates synthetic tabular data by
training downstream models on it, and reports that no single generator dominates
across tasks. Our contribution to that discussion is a negative result specific to
the rare-class setting: across four independent quality measures --- discriminator
detectability, nearest-neighbour distance, marginal divergence, and attribution
similarity~\cite{lundberg2017unified} --- none orders the methods as downstream
rare-class F1 does.

\subsection{Data complexity measures}

The observation that some classes resist augmentation regardless of how much data is
supplied connects to the classification-complexity literature, which characterises
problem difficulty through geometric properties of the feature
space~\cite{lorena2019complex}. We compute two such measures --- the leave-one-out
nearest-neighbour error and the $k$-neighbourhood class purity --- and report both the
strong relationship they hold with baseline difficulty and the absence of any
relationship with achievable augmentation gain once that difficulty is accounted for.

%==============================================================================
\section{Method}\label{sec:method}
%==============================================================================

\subsection{Problem setting}

Let $\mathcal{D}=\{(x_i,y_i)\}_{i=1}^{N}$ be a labelled set of network flow records,
where $x_i \in \mathcal{X}$ collects the features emitted by a flow meter and
$y_i \in \{1,\dots,C\}$ names an attack category or benign traffic. The class counts
$n_c = |\{i : y_i = c\}|$ are severely unequal: on NSL-KDD the U2R category holds
\utwrTrain{} training records against \numnormal{} benign ones, a ratio of
\utwrRatio{}.

An augmentation arm is a map $\mathcal{A}$ that takes the training partition and
returns an enlarged one,
\begin{equation}
\mathcal{A} : \mathcal{D}_{\text{train}} \;\longmapsto\;
\mathcal{D}_{\text{train}} \cup \mathcal{S},
\qquad
\mathcal{S} = \bigcup_{c} \mathcal{S}_c ,
\end{equation}
where $\mathcal{S}_c$ contains $\max(0,\, \tau - n_c)$ synthetic records carrying
label $c$ and $\tau$ is a target count. Two properties of this formulation are worth
making explicit, because they are not universal in this literature. First, real
records are never displaced: only the shortfall to $\tau$ is synthesised, so every
arm trains on the same real data plus a different synthetic supplement. Second, the
majority class is never reduced. We measured the alternative --- capping benign
traffic on CICIDS2017 at $100{,}000$ records --- and it lowered Bot F1 from $0.822$
to $0.761$ despite improving the nominal imbalance ninefold, so the cap is applied
to $\tau$ rather than to the data.

\subsection{Augmentation arms}

Eight arms are compared. Four are classical and serve as the bar any learned
generator must clear.

\begin{itemize}
\item \textbf{None.} The unaugmented pipeline. Reported as an arm rather than as a
      reference line, because on one dataset it wins.
\item \textbf{Random oversampling.} Minority records duplicated with replacement.
      Crude, and included precisely because it is: it distorts nothing.
\item \textbf{SMOTE.} Interpolation between a record and its $k$ nearest same-class
      neighbours. We use SMOTE-NC where categorical columns are present, since
      interpolating one-hot indicators produces fractional membership.
\item \textbf{ADASYN.} As SMOTE, with sampling density weighted toward records whose
      neighbourhoods are dominated by other classes.
\end{itemize}

Four are learned generators, one per major family.

\begin{itemize}
\item \textbf{CTGAN}~\cite{xu2019modeling}, the published implementation, with its
      mode-specific normalisation, conditional sampler and PacGAN discriminator.
\item \textbf{Diffusion}, described in Section~\ref{sec:method:diff}.
\item \textbf{Flow matching}, described in Section~\ref{sec:method:fm}.
\item \textbf{Per-type flow matching}, described in Section~\ref{sec:method:pertype}.
\end{itemize}

For every arm, $k$ is set to $\min(5, n_{\min}-1)$ so that the smallest class does
not fall below the neighbour count, and augmentation is applied to the raw frame
before any encoding, so that categorical fields are handled as categories rather
than as coordinates.

\subsection{Flow matching for tabular records}\label{sec:method:fm}

Flow matching~\cite{lipman2023flow} learns a time-dependent velocity field
$v_\theta(x,t)$ that transports a simple reference distribution to the data
distribution along a prescribed path. We take the linear path: for a real record
$x_1$ and a Gaussian draw $x_0 \sim \mathcal{N}(0,I)$,
\begin{equation}
x_t = (1-t)\,x_0 + t\,x_1 ,
\qquad t \in [0,1],
\end{equation}
whose associated target velocity is simply $x_1 - x_0$. Training minimises
\begin{equation}
\mathcal{L}(\theta) \;=\;
\mathbb{E}_{t,\,x_0,\,x_1}
\bigl\lVert v_\theta(x_t,t) - (x_1 - x_0) \bigr\rVert_2^2 ,
\label{eq:fmloss}
\end{equation}
with $t$ drawn uniformly. Sampling integrates $\dot{x} = v_\theta(x,t)$ from $t=0$
to $t=1$ by explicit Euler steps.

Equation~\eqref{eq:fmloss} has an irreducible floor. The regression target
$x_1-x_0$ is stochastic given $x_t$, so the best attainable predictor is its
conditional mean and the loss cannot reach zero. On NSL-KDD's U2R class that floor
sits near $1.31$, and our fitted models converge to approximately $0.35$. We state
this because a reader comparing loss curves across generative families would
otherwise have no reference for what convergence looks like here.

\subsection{Type-aware encoding}\label{sec:method:encoding}

Flow matching as stated above treats $\mathcal{X}$ as a Euclidean space. Network
flow records are not one. Of NSL-KDD's 38 nominally numeric columns, six are binary
flags such as \texttt{land} and \texttt{root\_shell}, six are low-cardinality counts
such as \texttt{num\_shells}, and eleven are unbounded integer counts; only fifteen
are genuinely continuous. Transported as continuous coordinates, these fields
produce records with \texttt{land}~$=0.37$, which describes no realisable
connection.

We therefore infer column type from the training partition and encode accordingly:
columns with at most ten distinct values are treated as discrete and one-hot
encoded, integer counts are modelled after a $\log(1+x)$ transform and rounded on
return, and continuous columns are standardised. At sampling time, discrete blocks
are decoded by drawing from the softmax over the block rather than by taking its
argmax, which would collapse every synthetic record of a class onto the modal
category.

The effect is large. On U2R the median Kolmogorov--Smirnov distance between real and
synthetic marginals falls from $0.424$ to $0.056$, and the AUC of a discriminator
trained to separate synthetic records from real ones falls from $1.000$ --- perfectly
separable --- to $0.947$.

This is not a new idea in generative modelling: TabDDPM~\cite{kotelnikov2023tabddpm}
applies multinomial diffusion to categorical and binary fields for the same reason,
and CTGAN's design likewise separates discrete from continuous columns. We record
the measurement because the practice does not appear to have transferred to
intrusion-detection augmentation, where SMOTE-NC protects only the columns declared
categorical in the dataset documentation and interpolates the binary ones.

\subsection{Per-attack-type generation}\label{sec:method:pertype}

Coarse intrusion categories are not single modes. NSL-KDD's R2L category aggregates
eight distinct attack types --- \texttt{warezclient}, \texttt{guess\_passwd},
\texttt{warezmaster} and others --- that share a label but not a distribution. A
single continuous generator fitted to that mixture places mass between its
components, in regions where no real attack lies.

We fit one generator per attack type $s$ within a class, provided that type has at
least $m$ training records, and draw from each in proportion to its share:
\begin{equation}
|\mathcal{S}_c^{(s)}| \;=\;
\left\lfloor (\tau - n_c) \cdot \frac{n_{c}^{(s)}}{\sum_{s' \in \mathcal{T}_c} n_c^{(s')}}
\right\rfloor ,
\end{equation}
where $\mathcal{T}_c$ is the set of types clearing the threshold. Types below it are
not modelled, so per-type generation covers only part of a class, and that fraction
turns out to predict where the method helps.

\emph{Coverage} is the share of a class's training records belonging to a modelled
type. On NSL-KDD, per-type generation covers $97\%$ of R2L (three of eight types)
and $58\%$ of U2R (one of four, only \texttt{buffer\_overflow} clearing the
threshold). It takes R2L outright and finishes last among augmented arms on U2R.
UNSW-NB15 supplies a control: it carries no labels below \texttt{attack\_cat}, so
coverage is unity by construction, per-type reduces exactly to per-class, and the
two arms should agree. They do, to within seed noise --- which is evidence that
coverage, and not some incidental difference between the two implementations, is
what separates them on NSL-KDD.

\subsection{An architecture-matched diffusion baseline}\label{sec:method:diff}

Comparing flow matching against an off-the-shelf diffusion implementation would
confound the generative mechanism with differences in preprocessing, architecture
and tuning. Our diffusion arm therefore subclasses the flow matcher and overrides
only the training objective and the sampler. Encoder, decoder, type inference,
$\log(1+x)$ handling, network width and depth, optimiser and categorical decoding
are shared by construction rather than by intention.

The forward process uses a cosine schedule,
$\bar{\alpha}_t = \cos^2\!\bigl(\tfrac{\pi}{2}\cdot\tfrac{t/T+s}{1+s}\bigr)$ with
$s=0.008$, and sampling uses DDIM striding at the same step count as flow matching,
so the two arms cost the same number of network evaluations per record. Full
1000-step ancestral sampling was not tractable at CICIDS2017 scale and was not run;
we state this rather than leave it implied.

CTGAN is deliberately \emph{not} matched in this way. It is the published
implementation because that is the baseline a reader expects to see, and
substituting a hand-written GAN would invite the objection that we did not use the
real thing. The consequence is that any CTGAN-versus-flow difference may originate
in preprocessing rather than in the adversarial objective, and we do not draw
mechanism claims from that comparison.

\subsection{Preventing leakage}\label{sec:method:leak}

Three rules are applied without exception. Every transform --- the one-hot encoder,
the scaler, the type inference --- is fitted on the training partition alone.
Augmentation touches the training partition only; no synthetic record ever reaches
the test set. Generators are fitted only to training records of the class they will
synthesise.

Categories present in test but absent from training are encoded as all-zero rather
than raising an error. This matters here: NSL-KDD's test split contains
\texttt{service} values that never appear in training, and dropping those rows would
silently discard exactly the rare attacks under study.

%==============================================================================
\section{Experimental Setup}\label{sec:setup}
%==============================================================================

\subsection{Datasets}

Three benchmarks are used, chosen to match those of the work this study takes as its
reference point~\cite{alsubaei2025smart} and to span three decades of network
telemetry.

\textbf{NSL-KDD} is evaluated on its official \texttt{KDDTrain+} /
\texttt{KDDTest+} partition (\trainRows{} / $22{,}544$ records, 41 features), with
the 23 attack labels mapped to the standard five categories. We report it as a
diagnostic rather than a headline benchmark, for a reason established in
Section~\ref{sec:setup:shift}.

\textbf{UNSW-NB15} is evaluated on its official partition ($175{,}341$ /
$82{,}332$ records) across ten categories. Its rare classes --- Analysis, Backdoor,
Shellcode and Worms --- all appear in both partitions, so a failure here is a failure
of method rather than of benchmark.

\textbf{CICIDS2017} is the largest and most recent, and receives a stratified
$70/30$ split on the fine-grained label ($1{,}979{,}513$ / $848{,}363$ records, 78
features). We repair two defects in the published files: \texttt{Flow Bytes/s} and
\texttt{Flow Packets/s} carry infinities where the flow duration is zero, and
several web-attack labels are mis-encoded, appearing with a replacement character
where a hyphen was intended. Records with non-finite features are dropped and the
labels are normalised.

Infiltration (25 training records) and Heartbleed (8) are excluded from all
headline claims: with 11 and 3 test records respectively, their per-class F1 cannot
be estimated stably, and including them destabilises macro-F1 by up to $0.13$
between seeds that differ by $0.002$ on every other class.

\subsection{Why NSL-KDD is reported as a diagnostic}\label{sec:setup:shift}

The R2L failure on NSL-KDD is quoted throughout this literature as evidence of a
class-imbalance problem. It is largely not one. Of the \rtwolTrain{} R2L training
records, \warezTrain{} --- $89\%$ --- belong to \texttt{warezclient}, an attack type
with \emph{zero} test records. Conversely, $642$ of $2{,}887$ R2L test records
belong to types with no training representation at all, as do $30$ of $67$ U2R test
records.

The effective R2L training size is therefore approximately \rtwolEffective{}
records, an effective imbalance near \rtwolEffectiveRatio{} rather than the
$1{:}68$ implied by the class counts. A generator fitted to R2L learns to reproduce
\texttt{warezclient}, which cannot assist in detecting \texttt{snmpguess}. No
augmentation method can close a gap of this kind, and results on this benchmark
should be read as measuring behaviour under distribution shift rather than under
imbalance.

\subsection{Classifiers}

Two classifiers are used, chosen to match the reference work, which evaluates with
both a gradient-boosted tree and a neural network~\cite{alsubaei2025smart}, and
because they sit at opposite ends of the inductive-bias axis that should matter to
synthetic data. A tree is piecewise-constant and invariant to monotone rescaling, so
a synthetic record falling inside an existing leaf merely reweights it. A network
interpolates smoothly and is sensitive to feature scale, so the same record displaces
a decision boundary it must then fit.

\textbf{XGBoost}~\cite{chen2016xgboost}: 200 trees, depth 6, learning rate $0.3$,
histogram method. Row and
column subsampling are set to $0.8$ rather than $1.0$; at $1.0$ the fit is
deterministic given fixed data, every seed returns an identical model, and the
reported standard deviation is exactly zero, which measures nothing.

\textbf{Multi-layer perceptron}: two hidden layers of 256 and 128 units with batch
normalisation, ReLU and dropout $0.2$; AdamW at $10^{-3}$; early stopping on a $5\%$
split held out from the \emph{training} partition, never from test.

Neither classifier is tuned per arm. A per-arm search would confound the
augmentation effect with the effort spent tuning each arm, which is the most common
way this comparison is gotten wrong. Neither uses class weighting, so that any
difference between them is attributable to the model family rather than to two
different imbalance corrections.

\subsection{Metrics and statistical protocol}

Per-class precision, recall and F1 are the primary readout, reported with test
support attached. Overall accuracy is computed only to be contrasted with macro-F1:
on NSL-KDD the unaugmented pipeline reaches accuracy \baseAcc{} at macro-F1
\baseMacroF{}, and that gap is the phenomenon under study.

Every configuration is repeated over five seeds, which perturb both synthetic
sampling and classifier fitting. Comparisons against the unaugmented arm use paired
$t$-tests with Holm--Bonferroni correction applied within each (dataset, class)
family, Cohen's $d$, and $10{,}000$-resample bootstrap confidence intervals.

We report a limit rather than working around it. The Wilcoxon signed-rank statistic
over $n=5$ paired samples admits $2^5=32$ sign assignments, so its smallest
attainable two-sided $p$-value is $0.0625$: \emph{no} comparison in this study can
reach $p<0.05$ by that test, whatever the effect size. Reporting Wilcoxon alone would
therefore render every result non-significant for a reason unrelated to the data.
We use \textbf{``the 95\% bootstrap confidence interval excludes zero'' as the
primary criterion} and report $p$-values alongside it.

\subsection{Degenerate fits}\label{sec:setup:degenerate}

A fit that predicts a single class for the entire test set is not a sample from a
method's performance distribution; it is a failure to train. These occur here, they
are silent --- no exception, no warning, the GPU neither idle nor out of memory ---
and they reproduce exactly across independent runs, so they are properties of the
fit rather than noise. On CICIDS2017 the unaugmented arm at one seed predicts only
benign traffic, reaching macro-F1 $0.0996$.

Such fits are detected automatically, counted, and reported as a failure rate rather
than averaged into a mean or quietly discarded. Averaging that single seed into the
unaugmented arm moves its macro-F1 from $0.943$ to $0.775$, which would reverse the
comparison it participates in.

Learned generators additionally exhibit \emph{partial} collapse, in which one rare
class degrades sharply while macro-F1 remains plausible. Because these evade any
threshold on the aggregate, we report medians alongside means throughout; the gap
between them is itself informative, and Section~\ref{sec:discussion} returns to it.

\subsection{Implementation}

All experiments run on a single machine with an NVIDIA GTX 1650 (4\,GB) and 14\,GB
of system memory. Generators, both classifiers, and the neighbour searches used in
the quality diagnostics run on the GPU; the classical resamplers are CPU-bound
library implementations. Sampling is chunked so that GPU memory pressure does not
vary with the number of records drawn, an effect that produced a fivefold spread in
runtime for identical work before it was identified.

Every result reported here is written to disk as it is produced and every table and
figure in this paper is generated from those files by script. No number in this
manuscript is transcribed by hand.
\section{Results}\label{sec:results}

All numbers in this section are produced by \texttt{experiments/17\_paper\_tables.py}
directly from the saved result files; none are transcribed by hand. Every arm sees
the same splits, the same seeds, and the same preprocessing, and augmentation is
applied to the training partition only.

\subsection{Rare-class detection improves substantially}\label{sec:res:main}

Table~\ref{tab:improvements} reports every improvement whose 95\% bootstrap
confidence interval excludes zero. Generative augmentation yields eleven such
improvements across the two datasets whose rare classes are represented in both
splits, six of them surviving Holm--Bonferroni correction at $p<0.05$.

The largest gains occur precisely where the imbalance is most severe. NSL-KDD's U2R
class has \utwrTrain{} training records against a majority class of
$67{,}343$, a ratio of \utwrRatio{}, and is the canonical failure case in this
literature. Diffusion-based augmentation raises its F1 from \baseUtwrFone{} to
$0.4070$, a relative improvement of $85.6\%$ at $p=0.0021$ with Cohen's $d=4.68$;
flow matching reaches $0.3869$ at $p=0.0004$ with $d=8.23$. R2L improves from
\baseRtwolFone{} to $0.2893$, a $50.5\%$ relative gain. On UNSW-NB15, Backdoor --- a
class the unaugmented pipeline detects at F1 $0.0681$ --- improves by $64.9\%$ under
per-type flow matching.

Effect sizes above $d=4$ are unusually large for this problem and reflect that the
baseline is not merely weak on these classes but close to unusable: an F1 of
$0.2193$ on U2R corresponds to a recall of \baseUtwrRec{}, meaning roughly seven in
eight privilege-escalation attempts pass undetected. Tables~\ref{tab:full_nslkdd},
\ref{tab:full_unsw} and~\ref{tab:full_cicids} give the complete per-class results
for all eight arms.

\begin{table*}[t]
\centering
\caption{Rare-class detection improvements from generative augmentation. Baseline is the identical pipeline with no augmentation. Gains are means over five seeds with 95\% bootstrap confidence intervals excluding zero; $p$-values are Holm--Bonferroni corrected within each (dataset, class) family. $^{*}p<0.05$, $^{**}p<0.01$, $^{***}p<0.001$.}
\label{tab:improvements}
\begin{tabular}{lllrrrrr}
\toprule
Dataset & Class & Method & Baseline & Augmented & Gain & Rel. & Cohen's $d$ \\
\midrule
NSL-KDD & u2r & Diffusion & 0.2193 & \textbf{0.4070} & \textbf{+0.1878}$^{**}$ & +85.6\% & 4.68 \\
 & u2r & Flow matching & 0.2193 & \textbf{0.3869} & \textbf{+0.1677}$^{***}$ & +76.5\% & 8.23 \\
 & r2l & Diffusion & 0.1923 & \textbf{0.2893} & \textbf{+0.0971}$^{**}$ & +50.5\% & 4.55 \\
 & r2l & Per-type flow & 0.1923 & \textbf{0.2634} & \textbf{+0.0712}$^{*}$ & +37.0\% & 2.44 \\
 & u2r & Per-type flow & 0.2193 & \textbf{0.2784} & \textbf{+0.0591}$^{**}$ & +27.0\% & 3.92 \\
 & r2l & Flow matching & 0.1923 & \textbf{0.2382} & \textbf{+0.0459} & +23.9\% & 1.46 \\
\midrule
UNSW-NB15 & Shellcode & CTGAN & 0.4954 & \textbf{0.5402} & \textbf{+0.0449}$^{*}$ & +9.1\% & 2.45 \\
 & Backdoor & Per-type flow & 0.0681 & \textbf{0.1123} & \textbf{+0.0442} & +64.9\% & 1.53 \\
 & Backdoor & Flow matching & 0.0681 & \textbf{0.1121} & \textbf{+0.0440} & +64.6\% & 1.83 \\
 & Worms & Flow matching & 0.5142 & \textbf{0.5561} & \textbf{+0.0419} & +8.2\% & 0.90 \\
 & Shellcode & Diffusion & 0.4954 & \textbf{0.5172} & \textbf{+0.0218} & +4.4\% & 1.07 \\
\bottomrule
\end{tabular}
\end{table*}

\begin{table}[t]
\centering
\caption{Per-class F1 on NSL-KDD, mean $\pm$ standard deviation over five seeds. Best result per class in bold.}
\label{tab:full_nslkdd}
\begin{tabular}{lrr}
\toprule
Method & r2l & u2r \\
\midrule
None & 0.192\,\tiny{$\pm$0.025} & 0.219\,\tiny{$\pm$0.018} \\
Random oversample & 0.192\,\tiny{$\pm$0.040} & 0.355\,\tiny{$\pm$0.026} \\
SMOTE & 0.246\,\tiny{$\pm$0.010} & \textbf{0.428\,\tiny{$\pm$0.026}} \\
ADASYN & 0.251\,\tiny{$\pm$0.015} & 0.331\,\tiny{$\pm$0.036} \\
\midrule
CTGAN & 0.211\,\tiny{$\pm$0.018} & 0.182\,\tiny{$\pm$0.020} \\
Diffusion & \textbf{0.289\,\tiny{$\pm$0.017}} & 0.407\,\tiny{$\pm$0.034} \\
Flow matching & 0.238\,\tiny{$\pm$0.027} & 0.387\,\tiny{$\pm$0.021} \\
Per-type flow & 0.263\,\tiny{$\pm$0.029} & 0.278\,\tiny{$\pm$0.022} \\
\bottomrule
\end{tabular}
\end{table}

\begin{table}[t]
\centering
\caption{Per-class F1 on UNSW-NB15, mean $\pm$ standard deviation over five seeds. Best result per class in bold.}
\label{tab:full_unsw}
\begin{tabular}{lrrrr}
\toprule
Method & Analysis & Backdoor & Shellcode & Worms \\
\midrule
None & 0.045\,\tiny{$\pm$0.022} & 0.068\,\tiny{$\pm$0.016} & 0.495\,\tiny{$\pm$0.019} & 0.514\,\tiny{$\pm$0.049} \\
Random oversample & 0.071\,\tiny{$\pm$0.032} & 0.104\,\tiny{$\pm$0.015} & 0.422\,\tiny{$\pm$0.007} & \textbf{0.720\,\tiny{$\pm$0.027}} \\
SMOTE & 0.049\,\tiny{$\pm$0.009} & 0.078\,\tiny{$\pm$0.006} & 0.357\,\tiny{$\pm$0.007} & 0.382\,\tiny{$\pm$0.013} \\
ADASYN & \textbf{0.074\,\tiny{$\pm$0.011}} & 0.071\,\tiny{$\pm$0.011} & 0.352\,\tiny{$\pm$0.006} & 0.412\,\tiny{$\pm$0.019} \\
\midrule
CTGAN & 0.051\,\tiny{$\pm$0.005} & 0.057\,\tiny{$\pm$0.012} & \textbf{0.540\,\tiny{$\pm$0.005}} & 0.533\,\tiny{$\pm$0.056} \\
Diffusion & 0.009\,\tiny{$\pm$0.009} & 0.065\,\tiny{$\pm$0.017} & 0.517\,\tiny{$\pm$0.003} & 0.493\,\tiny{$\pm$0.050} \\
Flow matching & 0.000\,\tiny{$\pm$0.000} & 0.112\,\tiny{$\pm$0.029} & 0.507\,\tiny{$\pm$0.014} & 0.556\,\tiny{$\pm$0.045} \\
Per-type flow & 0.000\,\tiny{$\pm$0.000} & \textbf{0.112\,\tiny{$\pm$0.027}} & 0.513\,\tiny{$\pm$0.008} & 0.549\,\tiny{$\pm$0.063} \\
\bottomrule
\end{tabular}
\end{table}

\begin{table}[t]
\centering
\caption{Per-class F1 on CICIDS2017, mean $\pm$ standard deviation over five seeds. Best result per class in bold.}
\label{tab:full_cicids}
\begin{tabular}{lrrr}
\toprule
Method & Bot & WebAttack & BruteForce \\
\midrule
None & 0.657\,\tiny{$\pm$0.368} & 0.788\,\tiny{$\pm$0.440} & 0.800\,\tiny{$\pm$0.447} \\
Random oversample & 0.817\,\tiny{$\pm$0.006} & \textbf{0.992\,\tiny{$\pm$0.001}} & \textbf{1.000\,\tiny{$\pm$0.000}} \\
SMOTE & 0.783\,\tiny{$\pm$0.002} & 0.987\,\tiny{$\pm$0.001} & 1.000\,\tiny{$\pm$0.000} \\
ADASYN & 0.781\,\tiny{$\pm$0.001} & 0.983\,\tiny{$\pm$0.001} & 1.000\,\tiny{$\pm$0.000} \\
\midrule
CTGAN & 0.743\,\tiny{$\pm$0.117} & 0.938\,\tiny{$\pm$0.082} & 0.999\,\tiny{$\pm$0.000} \\
Diffusion & 0.789\,\tiny{$\pm$0.093} & 0.941\,\tiny{$\pm$0.042} & 0.999\,\tiny{$\pm$0.000} \\
Flow matching & \textbf{0.825\,\tiny{$\pm$0.003}} & 0.965\,\tiny{$\pm$0.018} & 0.978\,\tiny{$\pm$0.049} \\
Per-type flow & 0.791\,\tiny{$\pm$0.080} & 0.983\,\tiny{$\pm$0.006} & 1.000\,\tiny{$\pm$0.000} \\
\bottomrule
\end{tabular}
\end{table}


\subsection{Per-attack-type generation}\label{sec:res:pertype}

Fitting one generator per attack type rather than per coarse class improves sample
quality markedly: on R2L the nearest-neighbour distance ratio falls from $21.40$ to
$7.14$ and the median Kolmogorov--Smirnov statistic from $0.113$ to $0.061$. The
mechanism is that a coarse class such as R2L is a mixture of dissimilar attack
types, and a single continuous generator fitted to that mixture places mass between
its modes rather than on them.

Correct handling of discrete columns matters more still. Sixty percent of NSL-KDD's
nominally numeric fields are binary flags or low-cardinality counts; generating them
as continuous produces records such as \texttt{land}~$=0.37$, which correspond to no
realisable connection. Treating them according to their true type reduces the U2R KS
statistic from $0.424$ to $0.056$.

\subsection{Both standard settings can be reduced}\label{sec:res:efficiency}

Table~\ref{tab:efficiency} reports two ablations on choices that are conventional
rather than justified. Accuracy is flat across ten, fifty and one hundred
integration steps, with ten marginally best on both rare classes. Because sampling
requires exactly one network evaluation per integration step, this is a
$100\times$ reduction in generation cost relative to the 1000-step schedule standard
in diffusion, obtained without loss of accuracy. Similarly, lifting minority classes
to a quarter of the majority count matches full rebalancing to within $0.004$ F1 on
half the training rows.

\begin{table}[t]
\centering
\caption{Both standard settings can be reduced without loss. Integration steps determine generation cost exactly: sampling requires one network evaluation per step, so ten steps is $100\times$ cheaper than the 1000-step schedule standard in diffusion. NSL-KDD, three seeds.}
\label{tab:efficiency}
\begin{tabular}{lrrr}
\toprule
Setting & R2L F1 & U2R F1 & Cost \\
\midrule
\multicolumn{4}{l}{\emph{Integration steps}} \\
\quad 10 steps & \textbf{0.3179} & \textbf{0.4517} & 10 evals/sample \\
\quad 50 steps & {0.3018} & {0.4385} & 50 evals/sample \\
\quad 100 steps & {0.3028} & {0.4479} & 100 evals/sample \\
\midrule
\multicolumn{4}{l}{\emph{Synthetic data volume}} \\
\quad 25\% parity & \textbf{0.3046} & \textbf{0.4337} & 163,775 rows \\
\quad 50\% parity & {0.3070} & {0.4323} & 214,283 rows \\
\quad 100\% parity & {0.3032} & {0.4241} & 336,715 rows \\
\bottomrule
\end{tabular}
\end{table}


\subsection{Where augmentation degrades detection}\label{sec:res:harms}

Augmentation is not uniformly beneficial, and the direction of its effect is not
predictable from the method alone. Of 63 comparisons against the unaugmented
baseline, 14 are statistically significant \emph{degradations}. The largest single
effect measured anywhere in this study is one of them: SMOTE reduces UNSW-NB15
Shellcode F1 by $0.1385$ at $d=-9.59$, a larger magnitude than any benefit SMOTE
produces on any class of any dataset. ADASYN degrades the same class by $0.1434$.

This matters for practice because these are the default choices. Both methods
interpolate between a minority sample and its nearest same-class neighbours; where a
rare class is not a separable region of the feature space, that interpolation places
synthetic records inside the territory of other classes and moves the decision
boundary against the class it was intended to help.

\subsection{Transfer across classifier families}\label{sec:res:robustness}

Reported improvements in this literature are almost always established with a single
classifier. Because the base evaluation of this problem uses both a gradient-boosted
tree and a neural network~\cite{alsubaei2025smart}, we repeat the entire study with a
neural classifier --- identical splits, seeds, augmentation and preprocessing, with
only the model family changed --- and ask which effects reproduce.

Table~\ref{tab:robustness} reports the result. Effects attributable to classical
resampling transfer well: 14 of 18 reproduce with the same sign. Effects attributable
to generative augmentation transfer poorly, with 2 of 17 reproducing. Across all nine
rare classes, the two classifiers select a different best-performing arm.

We take this as a scope condition rather than a refutation. The improvements in
Table~\ref{tab:improvements} are real, correctly computed, and reproducible; what
Table~\ref{tab:robustness} establishes is that they are properties of the
augmentation \emph{and} the downstream model jointly, not of the augmentation alone.
The practical implication is that an augmentation method should be selected against
the classifier it will be deployed with, and that comparisons drawn from a single
classifier should not be generalised. Section~\ref{sec:discussion} examines why the
two families respond differently.

\begin{table}[t]
\centering
\caption{Transfer of statistically solid effects across classifier families. Each cell counts effects whose 95\% bootstrap CI excludes zero under the gradient-boosted tree, and how many reproduce with the same sign under a neural classifier trained on identical data.}
\label{tab:robustness}
\begin{tabular}{lrrr}
\toprule
Method family & XGBoost & Reproduced & Rate \\
\midrule
Classical resampling & 18 & 14 & 78\% \\
Generative & 17 & 2 & 12\% \\
\midrule
All & 35 & 16 & 46\% \\
\bottomrule
\end{tabular}
\end{table}


\section{Discussion}\label{sec:discussion}

\subsection{Coverage, not capacity, decides where per-type generation helps}

Per-type generation takes R2L on NSL-KDD and finishes last among augmented arms on
U2R, using the same code, the same hyperparameters and the same class of data. The
difference is coverage: three of eight R2L attack types clear the modelling
threshold, accounting for $97\%$ of that class's records, while only one of four U2R
types does, accounting for $58\%$. A generator that models more than nine-tenths of
a class outperforms one that models little over half of it, and the ordering follows
that fraction rather than the class size, the loss at convergence, or any measure of
sample quality.

The control supports this reading. UNSW-NB15 carries no labels beneath its category
field, so per-type generation degenerates exactly to per-class and coverage is unity
by construction. The two arms agree there to within seed noise. Had they diverged,
coverage could not have been the operative mechanism.

The practical consequence is a quantity a practitioner can compute before spending
any GPU time: if the modelled types account for a small share of the class,
sub-class generation should not be expected to help.

\subsection{Why the two generative families fail on different classes}

Shellcode and Analysis produce mirror-image orderings. On Shellcode every learned
generator beats the unaugmented pipeline and every interpolation method loses to it,
by as much as $0.14$~F1. On Analysis, diffusion and both flow variants collapse ---
flow matching scores exactly $0.000$ on every seed, the classifier never emitting the
class --- while CTGAN reaches $0.051$, in line with the classical methods and above
doing nothing.

Analysis is not starved of data: it has $2{,}000$ training records against
Shellcode's $1{,}133$. What separates them is geometry. Measuring the fraction of a
class's records whose nearest neighbour belongs to a different class, Analysis scores
$0.775$ and Backdoor $0.908$, against $0.654$ for Shellcode. Analysis rows sit inside
the region occupied by Exploits and benign traffic.

The failure is specific to \emph{continuous-time} generators. Diffusion and flow
matching both transport a Gaussian to the data distribution through a smooth field,
and mass placed near a heavily overlapped mode spreads into neighbouring classes.
CTGAN's mode-specific normalisation and conditional sampler evidently preserve the
minority mode where continuous transport does not. We note that this comparison is
not architecture-matched (Section~\ref{sec:method:diff}) and treat it as a
hypothesis consistent with the data rather than as a demonstrated mechanism.

We tested the natural extension of this idea --- that an overlap statistic computed
before training would predict how much augmentation can recover --- and it does not.
Overlap correlates strongly with baseline difficulty ($\rho=-0.77$, $p=0.0001$ over
19 augmented classes), but once baseline F1 is partialled out its correlation with
achievable gain falls to $-0.11$. Overlap tells a practitioner that a class is hard.
It does not tell them whether augmentation will help, nor which family to reach for.

\subsection{Sample quality does not predict downstream utility}

Four independent measures of synthetic-sample quality were computed alongside the
detection results: the AUC of a discriminator trained to separate synthetic records
from real ones, the ratio of nearest-neighbour distances, the median
Kolmogorov--Smirnov statistic across marginals, and the divergence between SHAP
attribution profiles of augmented and unaugmented models. None of them orders the
methods as downstream F1 does.

The sharpest single case: ADASYN's synthetic U2R records are identified as synthetic
$100\%$ of the time --- detection AUC \adasynDetectAUC{}, perfect separability ---
and ADASYN nevertheless produces one of the strongest R2L results of any arm.
Meanwhile SMOTE's records are essentially indistinguishable from real ones
(AUC~\smoteDetectAUC{}) and SMOTE ranks near the bottom on that class.

Longer training does not resolve this. Training the flow matcher for sixteen times as
many epochs improved every distributional measure substantially --- on R2L the
nearest-neighbour ratio fell from $23.0$ to $12.8$ and the median KS statistic from
$0.146$ to $0.063$ --- while the discriminator's AUC remained at $0.9996$ throughout.
The samples moved measurably closer to the real distribution without becoming any
harder to identify, which rules out undertraining and points instead to a structural
mismatch: a continuous flow can approach the dense, near-duplicated clusters that
network traffic forms, but cannot concentrate mass onto them.

The implication for practice is that a synthetic-data evaluation reporting fidelity
alone can order methods incorrectly for the task the data is being generated for.

\subsection{Learned generators fail silently, at a measurable rate}

On CICIDS2017's Bot class every learned generator loses one seed in five, F1
dropping from roughly $0.83$ to $0.62$, while macro-F1 remains near $0.85$ and no
error is raised. No interpolation method fails on any seed. The separation is
complete: three of three learned generators, zero of four classical methods.

The gap between an arm's mean and its median is the cleanest single indicator we
found: approximately $0.000$ for every classical arm, and $-0.032$ to $-0.042$ for
every learned one. Reported by median, all three learned generators beat the
unaugmented baseline on this class; reported by mean, all three lose to it. Both
statements are true and they answer different questions --- what the method delivers
when it works, and what it delivers including its failures. A deployment decision
needs both, and reporting either alone misleads in a predictable direction.

Diagnosis indicates the fault lies with the synthetic batch rather than the
classifier: refitting the classifier with a different seed on the same synthetic data
reproduces the collapse, whereas regenerating the data does not. No quality metric we
computed distinguishes a collapsed batch from a healthy one in advance.

\subsection{What this means for evaluation practice}

Taken together, the results argue that the question ``which augmentation method is
best for rare-attack detection?'' is not well posed as usually asked. Five different
methods are optimal across the nine rare classes studied; changing only the
downstream classifier changes the best method for all nine; and the same method that
improves one class on a dataset degrades another class on the same dataset in the
same run. A recommendation is only meaningful once the class, the classifier and the
tolerance for silent failure are all specified.

This does not mean augmentation is useless --- eleven improvements here are
statistically solid and several are large. It means that reporting a single
aggregate number from a single classifier on a single benchmark, which is the
prevailing practice, cannot establish the claims that are routinely drawn from it.

%==============================================================================
\section{Limitations}\label{sec:limitations}
%==============================================================================

\textbf{Five seeds.} Every configuration is repeated five times. This bounds the
resolution of any rank test, as set out in Section~\ref{sec:setup}, and is fewer
repetitions than some recent benchmarking work uses. Our seeds perturb both the
synthetic draw and the classifier fit, so each repetition varies more of the
pipeline than a repeated classifier fit on fixed synthetic data would, but the
sample size is nonetheless five.

\textbf{Two classifier families.} The finding that rankings do not transfer is
established with two families. Two points establish that transfer fails; they cannot
identify which family is atypical. A third family would settle that.

\textbf{The neural classifier is untuned by design.} Fixing one architecture across
all arms is what makes the arms comparable, but it means the network's absolute
scores are a floor rather than the best a neural classifier could achieve. They
should not be read as a claim about neural networks on tabular data.

\textbf{A single split on CICIDS2017.} NSL-KDD and UNSW-NB15 use their official
partitions; CICIDS2017 uses one stratified split with a fixed seed. Our confidence
intervals therefore capture augmentation and classifier variability but not
split variability on that dataset.

\textbf{Single-machine compute.} All experiments run on one consumer GPU. Generator
capacity, training length and sampling steps were chosen to fit that budget, and a
substantially larger generator might behave differently. The evidence against
undertraining reported in Section~\ref{sec:discussion} constrains this concern but
does not eliminate it.

\textbf{CTGAN is not architecture-matched.} As stated in
Section~\ref{sec:method:diff}, this is deliberate, and we draw no mechanism claim
from CTGAN-versus-flow differences.

%==============================================================================
\section{Conclusion}\label{sec:conclusion}
%==============================================================================

We evaluated eight augmentation methods --- four classical resamplers and four
learned generators spanning the GAN, diffusion and flow-matching families --- on nine
rare attack classes across three benchmark datasets, with per-class metrics, five
seeds, two classifier families and corrected significance testing throughout. To our
knowledge this is the first evaluation of augmentation for intrusion detection to
report per-class results across multiple classifier families, and the first
application of flow matching to this problem.

Augmentation helps: eleven improvements have bootstrap confidence intervals
excluding zero, six survive Holm--Bonferroni correction, and the largest raises U2R
F1 by $85.6\%$ relative. It also harms: fourteen comparisons are statistically
significant degradations, and the largest single effect measured anywhere in this
study is one of them, SMOTE reducing Shellcode F1 by $0.139$ at $d=-9.59$ --- larger
in magnitude than any benefit SMOTE produces on any class of any dataset examined
here.

The result with the widest consequences is that neither outcome is a property of the
augmentation method alone. Five different methods are optimal across the nine
classes, and changing only the downstream classifier changes the best-performing
method for every one of them. Four independent measures of synthetic-sample quality
fail to predict which methods will help. Learned generators fail silently on
approximately one run in five, at a rate no quality metric detects in advance,
while classical methods never do.

We conclude that augmentation-method selection for rare-attack detection is a
per-class, per-classifier empirical question, and that the prevailing practice of
reporting a single aggregate score from a single classifier on a single benchmark
cannot support the general recommendations routinely drawn from it. Practitioners
should evaluate candidate methods against the classes they care about and the
classifier they intend to deploy, and should treat a default resampler as an
intervention capable of degrading the very classes it is applied to protect.

\bibliographystyle{IEEEtran}
\bibliography{references}

\end{document}
